\documentclass{article}

\input{../../Abgabenvorlage/pre.tex}

\newcommand{\blattnr}{1}
\newcommand{\ausgabe}{\dots}
\newcommand{\abgabe}{\dots}
\setarray{max. Punkte}{1;3;3;7}

\input{../../Abgabenvorlage/modul_infos.tex}
\input{../../Abgabenvorlage/deckblatt.tex}
\input{../../Abgabenvorlage/header.tex}
\input{../../Abgabenvorlage/mathmacros.tex}

\begin{document}

\exercise{1}

Hier Lösungen zu Aufgaben hinschreiben. Nach einen \verb!\exercise!-Befehl wird ein Zeilenumbruch erwartet.

\subexercise{1}
\blindtext[1]
\subexercise{3}
\blindtext[1]

\bigskip\medskip
\hrule

\bigexercise{1}

Falls die Aufgaben auf den Übungsblättern anders nummeriert sind, gibt es auch die folgende Darstellung für Aufgaben.

\bigsubexercise{1}
\begin{itemize*}
    \vspace*{-1.5em}
    \item[(i)] Erste Unterunteraufgabe
    \item[(ii)] Zweite Unterunteraufgabe
    \item[(iii)] Dritte Unterunteraufgabe
    \item[(iv)] Vierte Unterunteraufgabe
    % …
\end{itemize*}

\blindtext[1]
\bigsubexercise{3}
Sei $x\in\R$.

\zz.: Es gilt $x > 0$ oder $x\le0$
\begin{beweis}
Trivial.\done
\end{beweis}

\end{document}
